\section{Conclusion, Limitations, and Future Work}

We compared four regression models (MLP, DT, RF, mGBDT) on the \textit{USA House
Prices} dataset across four location-encoding variants, both as single models and as
a price / price-per-sqft stacking ensemble, all under shared 5-fold cross-validation.

\textbf{Key findings.}
(i) Coordinate and target encoding both beat naive categorical encoding for every
model, but their relative ranking is model-dependent and statistically inconclusive,
reflecting the irregular geography of housing demand.
(ii) Price-per-sqft stacking yields a further, consistent improvement for the stronger
models by decoupling price from living area and damping outliers.
(iii) The best overall configuration, an RF ensemble on coordinate-only features,
reaches a validation Adjusted $R^2$ of $0.794$ and cuts RMSE by $22\%$ relative to the
\texttt{rf}/\texttt{cat} baseline.
(iv) A multi-layer mGBDT did not help: its second layer fit noise, so a single
target-propagation layer (equivalent to plain XGBoost) was both simpler and stronger,
which we adopt as the final boosting model.

\textbf{Practical implication.}
The deeper, beyond-size relationships hypothesized at the outset did not emerge merely
from a more expressive model; they surfaced only after domain-aware preprocessing
(coordinate encoding and price-per-sqft stacking). Feature engineering, paired with a
well-regularized tree ensemble, was more decisive than model depth on this tabular task.

\textbf{Limitations.}
Our conclusions rest on a single dataset. The two techniques that drove the gains are
also dataset-specific: coordinate encoding presumes meaningful latitude/longitude
fields, and price-per-sqft stacking presumes a size-dominated, outlier-heavy target.
Both may not generalize to other markets or feature sets.

\textbf{Future work.}
Three directions follow naturally. First, validate coordinate encoding and
price-per-sqft stacking on additional datasets and housing markets to test how far the
gains transfer beyond Washington State. Second, instrument the pipeline to record the
secondary efficiency metrics (training time, inference time, peak memory), which the
current implementation does not yet measure. Third, revisit deeper mGBDT stacks with
stronger regularization or richer inverse mappings, to probe whether the layered
representation can help on tasks with more hierarchical structure than this one.
In future opportunities, we aim to apply mGBDT to complex datasets such as stock data, where more hierarchical learning may be possible.
